
\documentclass{article}
\usepackage{colortbl}
\usepackage{makecell}
\usepackage{multirow}
\usepackage{supertabular}

\begin{document}

\newcounter{utterance}

\centering \large Interaction Transcript for game `dond', experiment `semi\_en', episode 14 with reflexion{-}Qwen3.5{-}4B.
\vspace{24pt}

{ \footnotesize  \setcounter{utterance}{1}
\setlength{\tabcolsep}{0pt}
\begin{supertabular}{c@{$\;$}|p{.15\linewidth}@{}p{.15\linewidth}p{.15\linewidth}p{.15\linewidth}p{.15\linewidth}p{.15\linewidth}}
    \# & $\;$A & \multicolumn{4}{c}{Game Master} & $\;\:$B\\
    \hline

    \theutterance \stepcounter{utterance}  
    & & \multicolumn{4}{p{0.6\linewidth}}{
        \cellcolor[rgb]{0.9,0.9,0.9}{
            \makecell[{{p{\linewidth}}}]{
                \texttt{\tiny{[P1$\langle$GM]}}
                \texttt{You have attempted this same task before and did not succeed. These are the plans you wrote after those attempts.} \\
\texttt{Trial \#1: New plan: I failed because I treated the negotiation as a unilateral declaration rather than a two{-}way agreement. I sent a message stating my desired items ("1 pepper, 1 tomato") but did not explicitly confirm that the remaining items ("3 bags") were being assigned to the other player, nor did I ask for their confirmation. In this game, proposals must be complementary and mutually agreed upon; simply stating my preference without offering the specific counter{-}proposal for the other player's items (or explicitly stating the full split) likely caused the agreement to fail or the other player to reject it. Next time, I will immediately propose the complete division in a single message: "I take 1 pepper and 1 tomato; you take 3 bags." This clearly defines the complementary set for both parties, maximizing my score (7) while ensuring the other player receives a non{-}zero value (3), thereby satisfying the condition for a successful agreement.} \\
\\ 
\texttt{Trial \#2: New plan: I failed because I submitted a secret proposal prematurely without first establishing a verbal agreement with the other player. The rules state that proposals must be "complementary" and imply mutual acceptance; by sending the proposal immediately, I likely bypassed the negotiation phase where the other player would have confirmed or modified the split. Since I am the first to move, I should start with a clear, non{-}secret message explicitly stating the full division of items (e.g., "I take 1 pepper and 1 tomato; you take 3 bags") to invite their acceptance. Only after receiving their confirmation or if they agree verbally should I submit the secret proposal. This ensures the agreement is mutually recognized before the final step.} \\
\\ 
\texttt{Trial \#3: New plan: I failed because I treated the negotiation as a simple exchange of preferences rather than a binding agreement mechanism. In Trial \#1, I correctly identified the need to state the full split, but in Trial \#2, I realized that simply stating the split verbally is insufficient if the other player does not explicitly confirm it before I submit my secret proposal. The game requires "complementary" proposals to be accepted, which implies a mutual understanding. Since I am the first mover, I should not just state my take; I must explicitly offer the *entire* division and ask for the other player's confirmation of that specific split. If they agree verbally, I will then submit the secret proposal matching that agreed{-}upon split. If they do not agree or reject it, I will need to adjust my offer. Crucially, I must ensure my secret proposal exactly matches the verbal agreement to avoid the "zero points" penalty for non{-}complementary proposals. My next step is to send a message that says: "I propose you take 1 pepper and 1 tomato, and I take 3 bags. Please confirm this split." Only after receiving their confirmation will I submit the secret proposal: "{[}Proposal: 1 pepper, 1 tomato{]}".} \\
\\ 
\\ 
\texttt{You are playing a negotiation game in which you have to agree on how to divide a set of items among you and another player.} \\
\\ 
\texttt{Rules:} \\
\texttt{(a) You and the other player are given a set of items. Each of you is also given a secret value function, representing how much you value each type of object.} \\
\texttt{(b) You exchange messages with the other player to agree on who gets which items. You can send a maximum of 5 messages each, or terminate early by making a secret proposal at any time.} \\
\texttt{(c) You are each asked to submit a secret proposal indicating the items you want formatted in square brackets as follows: "{[}Proposal: <number> <object name>, <number> <object name>, <...>{]}"} \\
\texttt{(d) If your proposals are complementary, i.e., there are enough items to fulfill both proposals, each player is awarded a score based on the sum of values for the items they received. Otherwise, both of you get zero points.} \\
\texttt{(e) Your goal is to maximize the score you receive.} \\
\\ 
\texttt{Let us start.} \\
\\ 
\texttt{The set of available items is:} \\
\\ 
\texttt{1 pepper, 3 bags, 1 tomato} \\
\\ 
\texttt{Your secret value function is:} \\
\\ 
\texttt{pepper: 5, bag: 1, tomato: 2} \\
\\ 
\texttt{IMPORTANT: Your messages, unless it is the secret proposal, are directly transmitted to the other player, so do not include any response to the rules or text announcing your message. To make a secret proposal, use the indicated format. Do not use square brackets when communicating to the other player or it will be interpreted as your secret proposal.} \\
            }
        }
    }
    & \\ \\

    \theutterance \stepcounter{utterance}  
    & \multicolumn{4}{p{0.6\linewidth}}{
        \cellcolor[rgb]{0.9,0.9,0.9}{
            \makecell[{{p{\linewidth}}}]{
                \texttt{\tiny{[P1$\rangle$GM]}}
                \texttt{I propose you take 1 pepper and 1 tomato, and I take 3 bags. Please confirm this split.} \\
            }
        }
    }
    & & \\ \\

    \theutterance \stepcounter{utterance}  
    & & & \multicolumn{2}{p{0.3\linewidth}}{
        \cellcolor[rgb]{0.9,0.9,0.9}{
            \makecell[{{p{\linewidth}}}]{
                \texttt{\tiny{[GM$|$GM]}}
                \texttt{continue} \\
            }
        }
    }
    & & \\ \\

    \theutterance \stepcounter{utterance}  
    & & & \multicolumn{2}{p{0.3\linewidth}}{
        \cellcolor[rgb]{0.9,0.9,0.9}{
            \makecell[{{p{\linewidth}}}]{
                \texttt{\tiny{[GM$|$GM]}}
                \texttt{I propose you take 1 pepper and 1 tomato, and I take 3 bags. Please confirm this split.} \\
            }
        }
    }
    & & \\ \\

    \theutterance \stepcounter{utterance}  
    & & \multicolumn{4}{p{0.6\linewidth}}{
        \cellcolor[rgb]{0.9,0.9,0.9}{
            \makecell[{{p{\linewidth}}}]{
                \texttt{\tiny{[GM$\rangle$P2]}}
                \texttt{You have attempted this same task before and did not succeed. These are the plans you wrote after those attempts.} \\
\texttt{Trial \#1: New plan: I failed because I submitted a proposal for all items immediately, which likely resulted in a non{-}complementary agreement or a zero{-}score outcome since the other player could not fulfill my entire demand. In this negotiation game, I must first communicate with the other player to identify items they value highly and propose a split where we each get a subset of items that maximizes our individual scores. I should not submit a proposal for all items in the first turn. Instead, I will start by asking the other player for their proposal or suggesting a specific division (e.g., "I value bags most; do you want the pepper and tomato?") to ensure our proposals are complementary before finalizing the deal.} \\
\\ 
\texttt{Trial \#2: New plan: I failed because I submitted a proposal for all items in the first turn, which likely resulted in a non{-}complementary agreement or a zero{-}score outcome since the other player could not fulfill my entire demand. In this negotiation game, I must first communicate with the other player to identify items they value highly and propose a split where we each get a subset of items that maximizes our individual scores. I should not submit a proposal for all items in the first turn. Instead, I will start by asking the other player for their proposal or suggesting a specific division (e.g., "I value bags most; do you want the pepper and tomato?") to ensure our proposals are complementary before finalizing the deal.} \\
\\ 
\texttt{Trial \#3: New plan: I failed because I submitted a proposal for all items in the first turn, which likely resulted in a non{-}complementary agreement or a zero{-}score outcome since the other player could not fulfill my entire demand. In this negotiation game, I must first communicate with the other player to identify items they value highly and propose a split where we each get a subset of items that maximizes our individual scores. I should not submit a proposal for all items in the first turn. Instead, I will start by asking the other player for their proposal or suggesting a specific division (e.g., "I value bags most; do you want the pepper and tomato?") to ensure our proposals are complementary before finalizing the deal.} \\
\\ 
\\ 
\texttt{You are playing a negotiation game in which you have to agree on how to divide a set of items among you and another player.} \\
\\ 
\texttt{Rules:} \\
\texttt{(a) You and the other player are given a set of items. Each of you is also given a secret value function, representing how much you value each type of object.} \\
\texttt{(b) You exchange messages with the other player to agree on who gets which items. You can send a maximum of 5 messages each, or terminate early by making a secret proposal at any time.} \\
\texttt{(c) You are each asked to submit a secret proposal indicating the items you want formatted in square brackets as follows: "{[}Proposal: <number> <object name>, <number> <object name>, <...>{]}"} \\
\texttt{(d) If your proposals are complementary, i.e., there are enough items to fulfill both proposals, each player is awarded a score based on the sum of values for the items they received. Otherwise, both of you get zero points.} \\
\texttt{(e) Your goal is to maximize the score you receive.} \\
\\ 
\texttt{Let us start.} \\
\\ 
\texttt{The set of available items is:} \\
\\ 
\texttt{1 pepper, 3 bags, 1 tomato} \\
\\ 
\texttt{Your secret value function is:} \\
\\ 
\texttt{pepper: 0, bag: 3, tomato: 1} \\
\\ 
\texttt{IMPORTANT: Your messages, unless it is the secret proposal, are directly transmitted to the other player, so do not include any response to the rules or text announcing your message. To make a secret proposal, use the indicated format. Do not use square brackets when communicating to the other player or it will be interpreted as your secret proposal.} \\
\\ 
\\ 
\texttt{I propose you take 1 pepper and 1 tomato, and I take 3 bags. Please confirm this split.} \\
            }
        }
    }
    & \\ \\

    \theutterance \stepcounter{utterance}  
    & & & \multicolumn{4}{p{0.6\linewidth}}{
        \cellcolor[rgb]{0.9,0.9,0.9}{
            \makecell[{{p{\linewidth}}}]{
                \texttt{\tiny{[GM$\langle$P2]}}
                \texttt{I propose you take 1 pepper and 1 tomato, and I take 3 bags. Please confirm this split.} \\
            }
        }
    }
     \\ \\

    \theutterance \stepcounter{utterance}  
    & & & \multicolumn{2}{p{0.3\linewidth}}{
        \cellcolor[rgb]{0.9,0.9,0.9}{
            \makecell[{{p{\linewidth}}}]{
                \texttt{\tiny{[GM$|$GM]}}
                \texttt{continue} \\
            }
        }
    }
    & & \\ \\

    \theutterance \stepcounter{utterance}  
    & & & \multicolumn{2}{p{0.3\linewidth}}{
        \cellcolor[rgb]{0.9,0.9,0.9}{
            \makecell[{{p{\linewidth}}}]{
                \texttt{\tiny{[GM$|$GM]}}
                \texttt{I propose you take 1 pepper and 1 tomato, and I take 3 bags. Please confirm this split.} \\
            }
        }
    }
    & & \\ \\

    \theutterance \stepcounter{utterance}  
    & & \multicolumn{4}{p{0.6\linewidth}}{
        \cellcolor[rgb]{0.9,0.9,0.9}{
            \makecell[{{p{\linewidth}}}]{
                \texttt{\tiny{[P1$\langle$GM]}}
                \texttt{I propose you take 1 pepper and 1 tomato, and I take 3 bags. Please confirm this split.} \\
            }
        }
    }
    & \\ \\

    \theutterance \stepcounter{utterance}  
    & \multicolumn{4}{p{0.6\linewidth}}{
        \cellcolor[rgb]{0.9,0.9,0.9}{
            \makecell[{{p{\linewidth}}}]{
                \texttt{\tiny{[P1$\rangle$GM]}}
                \texttt{I take 3 bags.} \\
            }
        }
    }
    & & \\ \\

    \theutterance \stepcounter{utterance}  
    & & & \multicolumn{2}{p{0.3\linewidth}}{
        \cellcolor[rgb]{0.9,0.9,0.9}{
            \makecell[{{p{\linewidth}}}]{
                \texttt{\tiny{[GM$|$GM]}}
                \texttt{continue} \\
            }
        }
    }
    & & \\ \\

    \theutterance \stepcounter{utterance}  
    & & & \multicolumn{2}{p{0.3\linewidth}}{
        \cellcolor[rgb]{0.9,0.9,0.9}{
            \makecell[{{p{\linewidth}}}]{
                \texttt{\tiny{[GM$|$GM]}}
                \texttt{I take 3 bags.} \\
            }
        }
    }
    & & \\ \\

    \theutterance \stepcounter{utterance}  
    & & \multicolumn{4}{p{0.6\linewidth}}{
        \cellcolor[rgb]{0.9,0.9,0.9}{
            \makecell[{{p{\linewidth}}}]{
                \texttt{\tiny{[GM$\rangle$P2]}}
                \texttt{I take 3 bags.} \\
            }
        }
    }
    & \\ \\

    \theutterance \stepcounter{utterance}  
    & & & \multicolumn{4}{p{0.6\linewidth}}{
        \cellcolor[rgb]{0.9,0.9,0.9}{
            \makecell[{{p{\linewidth}}}]{
                \texttt{\tiny{[GM$\langle$P2]}}
                \texttt{{[}Proposal: 1 pepper, 1 tomato{]}} \\
            }
        }
    }
     \\ \\

    \theutterance \stepcounter{utterance}  
    & & & \multicolumn{2}{p{0.3\linewidth}}{
        \cellcolor[rgb]{0.9,0.9,0.9}{
            \makecell[{{p{\linewidth}}}]{
                \texttt{\tiny{[GM$|$GM]}}
                \texttt{proposal} \\
            }
        }
    }
    & & \\ \\

    \theutterance \stepcounter{utterance}  
    & & & \multicolumn{2}{p{0.3\linewidth}}{
        \cellcolor[rgb]{0.9,0.9,0.9}{
            \makecell[{{p{\linewidth}}}]{
                \texttt{\tiny{[GM$|$GM]}}
                [1, 0, 1]
            }
        }
    }
    & & \\ \\

    \theutterance \stepcounter{utterance}  
    & & \multicolumn{4}{p{0.6\linewidth}}{
        \cellcolor[rgb]{0.9,0.9,0.9}{
            \makecell[{{p{\linewidth}}}]{
                \texttt{\tiny{[P1$\langle$GM]}}
                \texttt{The other player has submitted a secret proposal. It is now your turn to submit a proposal of your own.} \\
            }
        }
    }
    & \\ \\

    \theutterance \stepcounter{utterance}  
    & \multicolumn{4}{p{0.6\linewidth}}{
        \cellcolor[rgb]{0.9,0.9,0.9}{
            \makecell[{{p{\linewidth}}}]{
                \texttt{\tiny{[P1$\rangle$GM]}}
                \texttt{{[}Proposal: 1 pepper, 1 tomato{]}} \\
            }
        }
    }
    & & \\ \\

    \theutterance \stepcounter{utterance}  
    & & & \multicolumn{2}{p{0.3\linewidth}}{
        \cellcolor[rgb]{0.9,0.9,0.9}{
            \makecell[{{p{\linewidth}}}]{
                \texttt{\tiny{[GM$|$GM]}}
                \texttt{proposal} \\
            }
        }
    }
    & & \\ \\

    \theutterance \stepcounter{utterance}  
    & & & \multicolumn{2}{p{0.3\linewidth}}{
        \cellcolor[rgb]{0.9,0.9,0.9}{
            \makecell[{{p{\linewidth}}}]{
                \texttt{\tiny{[GM$|$GM]}}
                [1, 0, 1]
            }
        }
    }
    & & \\ \\

    \theutterance \stepcounter{utterance}  
    & & & \multicolumn{2}{p{0.3\linewidth}}{
        \cellcolor[rgb]{0.9,0.9,0.9}{
            \makecell[{{p{\linewidth}}}]{
                \texttt{\tiny{[GM$|$GM]}}
                [[1, 0, 1], [1, 0, 1]]
            }
        }
    }
    & & \\ \\

    \theutterance \stepcounter{utterance}  
    & \multicolumn{4}{p{0.6\linewidth}}{
        \cellcolor[rgb]{0.95,0.95,0.95}{
            \makecell[{{p{\linewidth}}}]{
                \texttt{\tiny{[P2$\rightarrow$P1]}}
                \texttt{Post{-}episode Reflexion (not sent to GM):} \\
\texttt{New plan: I failed because I submitted a secret proposal before establishing a mutual verbal agreement, causing the game to treat my proposal as the final, unconfirmed offer rather than part of a negotiation. The rules require proposals to be "complementary" and accepted by both parties; by jumping straight to the secret format, I likely prevented the other player from confirming the split verbally, leading to a failure where the agreement was not mutually recognized. Next time, I will start by explicitly stating the full division of items ("I take 1 pepper and 1 tomato; you take 3 bags") and asking for their confirmation. Only after receiving their verbal agreement will I submit the secret proposal matching that exact split to ensure the conditions for a successful, complementary agreement are met.} \\
            }
        }
    }
    & & \\ \\

    \theutterance \stepcounter{utterance}  
    & & & \multicolumn{4}{p{0.6\linewidth}}{
        \cellcolor[rgb]{0.95,0.95,0.95}{
            \makecell[{{p{\linewidth}}}]{
                \texttt{\tiny{[P1$\rightarrow$P2]}}
                \texttt{Post{-}episode Reflexion (not sent to GM):} \\
\texttt{New plan: I failed because I submitted a secret proposal in the second turn while the game rules allow for up to 5 messages of negotiation before a proposal is required. By proposing immediately after the opponent's first move, I likely forced a non{-}complementary split or missed an opportunity to negotiate a better distribution. The correct strategy is to use the first message to propose a specific division (e.g., "I take the pepper and tomato, you take the bags") to see if the opponent agrees or counters. If they agree, we can terminate early. If they disagree, I will use subsequent messages to adjust the split based on their counter{-}offer, ensuring that by the time I submit a formal secret proposal, the items are fully allocated and complementary to maximize my score.} \\
            }
        }
    }
     \\ \\

\end{supertabular}
}

\end{document}
